\section{Πρόβλημα 2}

Τα παρακάτω αφορούν δικτυακά παίγνια συμφόρησης με 2 (ατομικούς, ισοβαρείς)
παίκτες και γραμμικές συναρτήσεις καθυστέρησης.

\subsection{Ερώτηση (i)}

Θα βρούμε παίγνιο όπου οι δύο παίκτες μοιράζονται κοινή αφετηρία αλλά έχουν
πιθανώς διαφορετικό τερματισμό και το τίμημα της αναρχίας είναι 2.

Θεωρούμε το ακόλουθο δίκτυο, όπου ο παίκτης 1 θέλει να πάει από $s$ στο $t_1$
και ο παίκτης 2 από $s$ στο $t_2$:

\begin{center}
\begin{tikzpicture}[
  node distance=1.5cm and 2.2cm,
  every node/.style={circle, draw, fill=yellow!70!orange, minimum size=0.85cm, font=\bfseries},
  >={Stealth[length=6pt]}
]
  \node (s)  {$s$};
  \node (a)  [above right=of s] {$a$};
  \node (b)  [below right=of s] {$b$};
  \node (t1) [right=of a] {$t_1$};
  \node (t2) [right=of b] {$t_2$};

  \draw[->] (s) -- node[above left, draw=none, fill=none, font=\normalfont] {$x$} (a);
  \draw[->] (s) -- node[below left, draw=none, fill=none, font=\normalfont] {$x$} (b);
  \draw[->] (a) -- node[above,      draw=none, fill=none, font=\normalfont] {$0$} (t1);
  \draw[->] (b) -- node[below,      draw=none, fill=none, font=\normalfont] {$0$} (t2);
  \draw[->] (a) to[bend right=0] node[right, draw=none, fill=none, font=\normalfont] {$1$} (t2);
  \draw[->] (b) to[bend left=0]  node[left,  draw=none, fill=none, font=\normalfont] {$1$} (t1);
\end{tikzpicture}
\end{center}

Ο Κάθε παίκτης έχει δύο πιθανά μονοπάτια. Ο παίκτης 1 μπορεί να πάει είτε
$s \to a \to t_1$ (κόστος $\ell_{sa}(x) + 0$) είτε $s \to b \to t_1$ (κόστος
$\ell_{sb}(x) + 1$). Συμμετρικά ο παίκτης 2 μπορεί να πάει είτε
$s \to b \to t_2$ είτε $s \to a \to t_2$.

Θεωρούμε το προφίλ όπου ο παίκτης 1 επιλέγει $s \to b \to t_1$ και ο παίκτης 2
επιλέγει $s \to a \to t_2$. Τότε αυτό αποτελεί σημείο ισορροπίας:
\begin{itemize}
  \item Παίκτης 1: $\ell_{s \to b}(1) + 1 = 2$. Αν αλλάξει σε $s \to a \to t_1$
    θα έχει $\ell_{s \to a}(2) + 0 = 2$.
  \item Παίκτης 2: $\ell_{s \to a}(1) + 1 = 2$. Αν αλλάξει σε $s \to b \to t_2$
    θα έχει $\ell_{s \to b}(2) + 0 = 2$.
\end{itemize}
Το κοινωνικό κόστος είναι $2 + 2 = 4$.

Θεωρούμε το προφίλ όπου ο παίκτης 1 επιλέγει $s \to a \to t_1$ και ο παίκτης 2
επιλέγει $s \to b \to t_2$.
\begin{itemize}
  \item Παίκτης 1: $\ell_{s \to a}(1) + 0 = 1$.
  \item Παίκτης 2: $\ell_{s \to b}(1) + 0 = 1$.
\end{itemize}
Κοινωνικό κόστος $1 + 1 = 2$. Άρα το τίμημα της αναρχίας είναι
$\frac{4}{2} = 2$.

\subsection{Ερώτηση (ii)}

Θα δείξουμε ότι το τίμημα της αναρχίας δεν μπορεί να είναι χειρότερο από 2.

Έστω $f$ μια ισορροπία Nash και $f^*$ η κοινωνικά βέλτιστη ροή σε ένα δικτυακό
παιχνίδι συμφόρησης με 2 ατομικούς, ισοβαρείς παίκτες, κοινή αφετηρία $s$ και
(πιθανώς διαφορετικούς) τερματισμούς $t_1, t_2$. Έστω επίσης $p_1, p_2$ και
$p_1^*, p_2^*$ τα αντίστοιχα προφίλ στρατηγικών κάθε παίκτη. Κάθε ακμή έχει
γραμμική συνάρτηση καθυστέρησης $c_e(x) = a_e x + b_e$ με $a_e, b_e \geq 0$.
Συμβολίζουμε με $C(f) = \sum_{e \in E} f_e \cdot c_e(f_e)$ το κοινωνικό κόστος.
Συμβολίζουμε με $C_1, C_2$ τα κόστη στο προφίλ $(p_1, p_2)$ για τον κάθε παίκτη
αντίστοιχα, και $C_1^*, C_2^*$ τα κόστη του κάθε παίκτη στα προφίλ
$(p_1^*, p_2)$ και $(p_1, p_2^*)$ αντίστοιχα. Τέλος, για κάθε ακμή $e$ ορίζουμε
τους δείκτες $p^i(e) = \mathbf{1}[e \in p_i]$, $q^i(e) = \mathbf{1}[e \in p_i^*]$
και τις ροές $f_e = p^1(e) + p^2(e)$, $f_e^* = q^1(e) + q^2(e)$.

Από τον ορισμό του σημείου ισορροπίας, κανένας παίκτης δεν βελτιώνεται
αλλάζοντας μονοπάτι, άρα
\begin{equation}\label{eq.nash}
  C_1 \leq C_1^* \text{ και } C_2 \leq C_2^* \,.
\end{equation}
Όταν ο παίκτης $i$ αλλάζει στο βέλτιστο μονοπάτι του $p_i^*$ ενώ ο άλλος
παραμένει στο $p_{-i}$, για κάθε ακμή $e \in p_i^*$ η ροή γίνεται
\begin{equation}\label{eq.flow}
  f_{e, i}^* = 1 + p^{-i}(e) \,,
\end{equation}
αφού ο σίγουρα θα την χρησιμοποιεί πλέον ο παίκτης $i$, ενώ ο άλλος παίκτης δεν
αλλάζει στρατηγική. Οπότε
$C_i^* = \sum_{e \in p_i^*} c_e(1 + p^{-i}(e))$. Αθροίζοντας την \eqref{eq.nash}
για $i = 1, 2$:
\begin{equation}\label{eq.sum}
  C(f) = C_1 + C_2 \leq C_1^* + C_2^* = \sum_{e \in p_1^*} c_e(1 + p^2(e))
    + \sum_{e \in p_2^*} c_e(1 + p^1(e)) \,.
\end{equation}

Χρησιμοποιούμε τους όρους $q^i(e)$ για να ενώσουμε τα δύο αθροίσματα. Κάθε
άθροισμα γίνεται
$$
  \sum_{e \in p_i^*} c_e(1 + p^{-i}(e)) = \sum_{e \in E} q^i(e) \cdot c_e(1 + p^{-i}(e)) \,.
$$
Και συνολικά η \eqref{eq.sum} γράφεται ως
\begin{equation}\label{eq.edges}
  C(f) \leq \sum_{e \in E} \bigl( q^1(e) \cdot c_e(1 + p^2(e)) + q^2(e) \cdot c_e(1 + p^1(e)) \bigr)\,.
\end{equation}

Θα δείξουμε ότι
$$
  \sum_{e \in E} \bigl( q^1(e) \cdot c_e(1 + p^2(e)) + q^2(e) \cdot c_e(1 + p^1(e)) \bigr)\,.
  \leq 2C(f^*) = 2 \sum_{e \in E} f_e^* \cdot c_e(f_e^*) \,,
$$
δείχνοντας ότι κάθε στοιχείο του αριστερά μέλους είναι μικρότερο από το
αντίστοιχο δεξιά. Για συντομία, ορίζουμε
$X_e = q^1(e) \cdot c_e(1 + p^2(e)) + q^2(e) \cdot c_e(1 + p^1(e))$.
Άρα, θα δείξουμε ότι για κάθε ακμή $e \in E$ ισχύει
$X_e \leq 2\, f_e^* \cdot c_e(f_e^*)$, διακρίνοντας περιπτώσεις στη βέλτιστη
ροή $f_e^* = q_1(e) + q_2(e) \in \{0, 1, 2\}$.

\begin{itemize}
  \item $f_e^* = 0$: τότε $q_1(e) = q_2(e) = 0$, άρα $X_e = 0 \leq 0$.

  \item $f_e^* = 1$: χωρίς βλάβη της γενικότητας $q_1(e) = 1, q_2(e) = 0$. Τότε
    $
      X_e = c_e(1 + p_2(e)) \leq c_e(2) = 2 a_e + b_e \,,
    $
    ενώ $2\, f_e^*\, c_e(f_e^*) = 2\, c_e(1) = 2 a_e + 2 b_e$. Αφού $b_e \geq 0$,
    έχουμε $X_e \leq 2\, c_e(1)$.

  \item $f_e^* = 2$: $q_1(e) = q_2(e) = 1$, άρα
    $
      X_e = c_e(1 + p_2(e)) + c_e(1 + p_1(e)) \leq 2\, c_e(2) = 4 a_e + 2 b_e \,,
    $
    ενώ $2\, f_e^*\, c_e(f_e^*) = 4\, c_e(2) = 8 a_e + 4 b_e$.
\end{itemize}

Άρα,
$$
  C(f) \leq \sum_{e \in E} X_e \leq 2 \sum_{e \in E} f_e^* \cdot c_e(f_e^*) = 2\, C(f^*) \,.
$$
και το τίμημα της αναρχίας είναι $C(f) / C(f^*) \leq 2$.

\subsection{Ερώτηση (iii)}

To παίγνιο του παρακάτω σχήματος για δύο παίκτες με κοινή αφετηρία το $s$ και
τερματισμό $t$ έχει τίμημα της αναρχίας $8/5$.

Παρότι το κάτω μονοπάτι έχει $b = -2$, στην περίπτωση ατομικού παιγνίου μια ακμή
που χρησιμοποιείται έχει ροή τουλάχιστον $1$, οπότε το κόστος που βλέπει ένας
παίκτης δεν είναι ποτέ αρνητικό, αλλά τουλάχιστον $3 \cdot 1 - 2 = 1$.

\begin{center}
\begin{tikzpicture}[
  node distance=2.5cm,
  every node/.style={circle, draw, fill=yellow!70!orange, minimum size=0.85cm, font=\bfseries},
  >={Stealth[length=6pt]}
]
  \node (s) {$s$};
  \node (t) [right=of s] {$t$};

  \draw[->] (s) to[bend left=35]  node[above, draw=none, fill=none, font=\normalfont] {$3x + 1$} (t);
  \draw[->] (s) to[bend right=35] node[below, draw=none, fill=none, font=\normalfont] {$3x - 2$} (t);
\end{tikzpicture}
\end{center}


Τα δυνατά προφίλ στρατηγικών είναι:
\begin{itemize}
  \item Και οι δύο παίκτες στο πάνω μονοπάτι: κάθε παίκτης πληρώνει
    $3 \cdot 2 + 1 = 7$. Δεν είναι σημείο ισορροπίας γιατί αν κάποιος αλλάξει
    στο κάτω μονοπάτι θα πληρώσει $3 \cdot 1 - 2 = 1 < 7$. Κοινωνικό κόστος
    $7 + 7 = 14$.
  \item Ένας παίκτης στο πάνω και ένας στο κάτω: ο πάνω πληρώνει
    $3 \cdot 1 + 1 = 4$ και ο κάτω πληρώνει $3 \cdot 1 - 2 = 1$. Είναι σημείο
    ισορροπίας:
    \begin{itemize}
      \item Ο πάνω αν αλλάξει στο κάτω θα πληρώσει $3 \cdot 2 - 2 = 4$, ίδιο
        με τώρα.
      \item Ο κάτω αν αλλάξει στο πάνω θα πληρώσει $3 \cdot 2 + 1 = 7 > 1$.
    \end{itemize}
    Κοινωνικό κόστος $4 + 1 = 5$.
  \item Και οι δύο παίκτες στο κάτω μονοπάτι: κάθε παίκτης πληρώνει
    $3 \cdot 2 - 2 = 4$. Είναι σημείο ισορροπίας γιατί αν κάποιος αλλάξει στο
    πάνω θα πληρώσει $3 \cdot 1 + 1 = 4$, ίδιο με τώρα. Κοινωνικό κόστος
    $4 + 4 = 8$.
\end{itemize}

Από τα παραπάνω προφίλ, δύο είναι σημεία ισορροπίας με κοινωνικά κόστη $5$ και
$8$, οπότε το χειρότερο σημείο ισορροπίας έχει κόστος $8$ (και οι δύο παίκτες
στο κάτω μονοπάτι). Το βέλτιστο κοινωνικό κόστος είναι το ελάχιστο από όλα τα
προφίλ ($14$, $5$, $8$), δηλαδή $5$ (ένας παίκτης σε κάθε μονοπάτι). Άρα το
τίμημα της αναρχίας είναι $\frac{8}{5}$.
